\section{Конъюнктивные и булевы языки в качестве ограничений на пути}
\label{sec:ConjunctiveBooleanConstraints}
\tikzsetfigurename{FLPQ_ConjBoolConstr_}

Использование многокомпонентных контекстно-свободных языков, рассмотренных в предыдущем разделе, позволяет описывать ограничения на пути, не выразимые в терминах обычных КС-языков.
Однако и этот класс имеет свои пределы.
В частности, MCFL не замкнуты относительно операции пересечения, что не позволяет одновременно отслеживать несколько независимых сбалансированных свойств пути с гарантированной точностью.
Конъюнктивные и булевы грамматики, введённые в главе~\ref{chpt:ConjunctiveBoolean}, снимают это ограничение: операция конъюнкции встроена непосредственно в механизм вывода, что позволяет в одной грамматике комбинировать несколько условий на структуру пути.

На практике такая выразительность востребована, например, в статическом анализе программ.
Точный \emph{context-sensitive} и \emph{data-dependence} анализ требует одновременного отслеживания сбалансированности вызовов процедур и обращений к полям/указателям, что моделируется шафлом языков Дика и не является контекстно-свободным~\sidecite{reps2000undecidability}.
При использовании КС-ограничений одно из свойств приходится аппроксимировать регулярным языком~\sidecite{huang2015scalable,sridharan2006refinement}, теряя точность.
\emph{Линейные конъюнктивные грамматики} (LCL), напротив, позволяют сохранить оба свойства точными и демонстрируют значительные преимущества в точности и масштабируемости~\sidecite{zhang2017context}.

Перейдём к уточнению формальной постановки задачи для конъюнктивных и булевых ограничений.
Пусть дан помеченный граф $\mathcal{G} = \langle V, E, \Sigma \rangle$ и конъюнктивная (соответственно, булева) грамматика $G = \langle \Sigma, N, P, S \rangle$.
В \emph{реляционной семантике} запросов для каждого нетерминала $A \in N$ определяется отношение
\[
    R_A = \{(i,j) \in V \times V \mid \exists \pi: i \pi j \land l(\pi) \in L(A)\},
\]
и результатом запроса для стартового нетерминала $S$ является множество пар вершин $R_S$.
Задача состоит в нахождении (или аппроксимации) всех таких пар.

Ключевое отличие от случая КС- и MCFG-ограничений~--- вопрос разрешимости.
Известно, что задача поиска путей с ограничениями в виде конъюнктивных языков в реляционной семантике в общем случае \emph{неразрешима}~\sidecite{!!!}.
Это фундаментальное ограничение вынуждает использовать аппроксимационный подход: строить \emph{аппроксимацию сверху}~--- множество пар вершин, которое гарантированно содержит все пары, действительно связанные путём с требуемым ограничением, но может включать и <<лишние>> пары.
То есть алгоритм обладает свойством полноты (все действительные пути будут найдены), но не обязательно корректен (могут быть ложноположительные ответы).

Конкретные алгоритмы, реализующие такой аппроксимационный подход,~--- матричный алгоритм для конъюнктивных грамматик и его расширение на случай булевых грамматик над ациклическими графами~--- подробно рассматриваются в главе~\ref{chpt:BeyondCFL}.
Общая идея состоит в том, чтобы определить конъюнктивное произведение матриц, обобщающее обычное матричное умножение над полукольцом на случай нескольких конъюнктов в правой части правила, и итеративно вычислять транзитивное замыкание до достижения неподвижной точки.
